%% =====================================================================
%%  B1-T-03-02 -- what B1 contributes to "The Operator Profile"
%%  -------------------------------------------------------------------
%%  LAYER A: APPEND-ONLY. Written once, then left alone. Material found
%%  only here sits in the `unique` slot and is protected by rule R10.
%% =====================================================================
%% @kind: source
%% @id: B1-T-03-02
%% @book: B1
%% @topic: T-03-02
%% @locator: Tactic 1, pp. 9--18
%% @status: distilled
%% @unique: yes
%% @integrated: yes
%% @updated: 2026-09-19

\dossier{B1}{T-03-02}{\bookshortB1}{Tactic 1, pp. 9--18}

\begin{distilled}
  \item Character types to watch out for, drawn from the author's direct experience of being defrauded repeatedly by people he liked.
  \item The over-reassurer: every con artist took pains to assert his trustworthiness before the loss, using variations on the same formula --- you will not regret this, I have always paid you, you can count on me.
  \item A running list of the author's own losses is given with the exact reassurance used in each case, which is the evidentiary basis of the type.
  \item The once-wealthy man: nearly every operator had a story of money formerly held and lost.
  \item Demonstrative gratitude and physical warmth --- back-slapping, repeated assurance that you will get everything coming to you --- are described as the position from which the knife goes in.
  \item Terms that favour the target are said to speak for themselves; a person who repeats your benefit at length is selling something.
\end{distilled}

\begin{unique}
  \item The over-reassurance pattern documented against the author's own recorded losses, with the operator's exact formulae, and the derived rule that favourable terms do not need to be argued.
  \item The once-wealthy story as a distinct type marker rather than as background detail.
\end{unique}

\begin{terms}
  \item \emph{protest too much} --- asserting a quality so insistently that the assertion itself becomes the evidence against it.
\end{terms}
